% GENERATED FILE. Edit canonical lessons, then run npm run generate:books.
% canonical-source-hash: fnv1a64:cf0bb198c1d66eca
% canonical-lessons: GU-C365-siddhat, GU-C365-vyavhar, GU-C365-anuman, GU-C365-vastavik, GU-C365-sakalvu

\chapter{Theory and Practice}
\label{ch:gu-theory-and-practice}
\hlchaptermodality{365}

\begin{chapteropening}
\textbf{By the end of this chapter:} \emph{I can say whether an idea works in practice.}

Complete the last lesson of chapter 365: I can say whether an idea works in practice.
\end{chapteropening}

\section[\texorpdfstring{\gu{સિદ્ધાંત} (siddhā̃t)}{સિદ્ધાંત (siddhā̃t)}]{\gu{સિદ્ધાંત} (siddhā̃t) — a principle (a theory)}

\label{lesson:GU-C365-siddhat}



\begin{quote}
Before the new one: say the Gujarati for a mystery, then the Gujarati for to uncover.
\end{quote}

\subsection*{You'll want to know: \gu{સિદ્ધાંત}}
\textbf{\gu{સિદ્ધાંત}} — \emph{siddhā̃t} — \textquotedblleft{}a principle (a theory)\textquotedblright{}.

\begin{quote}
\gu{આ} \gu{સિદ્ધાંત} \gu{કાગળ} \gu{પર} \gu{સારો} \gu{છે}.
\end{quote}

\begin{quote}
\emph{ā siddhā̃t kāgaḷ par sāro che.}
\end{quote}

\begin{quote}
This theory looks good on paper.
\end{quote}

A masculine noun.

\begin{cousinweb}
\textbf{\gu{સિદ્ધાંત}} (\emph{siddhā̃t}) is a tatsama, borrowed from Sanskrit \emph{siddhānta}, \textquotedblleft{}an established end, a settled doctrine\textquotedblright{} (Monier-Williams).
\end{cousinweb}

\begin{grammarlens}[title={What you've built}]
One more word for giving reasons and opinions.
\end{grammarlens}

\subsection*{Your turn}
\begin{itemize}
\raggedright
  \item \emph{Say it:} \emph{siddhā̃t}
  \item \emph{Say it:} \emph{siddhā̃t}, once more
  \item \emph{Say it:} \emph{ā siddhā̃t kāgaḷ par sāro che.} — This theory looks good on paper.
  \item \emph{Recall:} say the Gujarati for a mystery, then the Gujarati for to uncover, then say \emph{siddhā̃t} again
\end{itemize}

\begin{tcolorbox}[breakable,colback=teal!4,colframe=teal!35!black,title={Before you move on}]
What does \gu{સિદ્ધાંત} mean? (\textquotedblleft{}A principle (a theory)\textquotedblright{}.) Say it once more.
\end{tcolorbox}
\section[\texorpdfstring{\gu{વ્યવહાર} (vyavhār)}{વ્યવહાર (vyavhār)}]{\gu{વ્યવહાર} (vyavhār) — everyday practice; dealings}

\label{lesson:GU-C365-vyavhar}



\begin{quote}
Before the new one: say the Gujarati for to uncover, then the Gujarati for a principle.
\end{quote}

\subsection*{You'll want to know: \gu{વ્યવહાર}}
\textbf{\gu{વ્યવહાર}} — \emph{vyavhār} — \textquotedblleft{}everyday practice; dealings\textquotedblright{}.

\begin{quote}
\gu{પણ} \gu{વ્યવહારમાં} \gu{આ} \gu{સિદ્ધાંત} \gu{ચાલતો} \gu{નથી}.
\end{quote}

\begin{quote}
\emph{paṇ vyavhārmā̃ ā siddhā̃t cālto nathī.}
\end{quote}

\begin{quote}
But in practice this theory does not work.
\end{quote}

A masculine noun.

\begin{cousinweb}
\textbf{\gu{વ્યવહાર}} (\emph{vyavhār}) is a tatsama, borrowed from Sanskrit \emph{vyavahāra}, \textquotedblleft{}doing, practice, conduct\textquotedblright{} (Monier-Williams).
\end{cousinweb}

\begin{grammarlens}[title={What you've built}]
One more word for giving reasons and opinions.
\end{grammarlens}

\subsection*{Your turn}
\begin{itemize}
\raggedright
  \item \emph{Say it:} \emph{vyavhār}
  \item \emph{Say it:} \emph{vyavhār}, once more
  \item \emph{Say it:} \emph{paṇ vyavhārmā̃ ā siddhā̃t cālto nathī.} — But in practice this theory does not work.
  \item \emph{Recall:} say the Gujarati for to uncover, then the Gujarati for a principle, then say \emph{vyavhār} again
\end{itemize}

\begin{tcolorbox}[breakable,colback=teal!4,colframe=teal!35!black,title={Before you move on}]
What does \gu{વ્યવહાર} mean? (\textquotedblleft{}Everyday practice; dealings\textquotedblright{}.) Say it once more.
\end{tcolorbox}
\section[\texorpdfstring{\gu{અનુમાન} (anumān)}{અનુમાન (anumān)}]{\gu{અનુમાન} (anumān) — an inference}

\label{lesson:GU-C365-anuman}



\begin{quote}
Before the new one: say the Gujarati for a principle, then the Gujarati for everyday practice.
\end{quote}

\subsection*{You'll want to know: \gu{અનુમાન}}
\textbf{\gu{અનુમાન}} — \emph{anumān} — \textquotedblleft{}an inference\textquotedblright{}.

\begin{quote}
\gu{આકાશ} \gu{જોઈને} \gu{દાદાએ} \gu{વરસાદનું} \gu{અનુમાન} \gu{કર્યું}.
\end{quote}

\begin{quote}
\emph{ākāś joīne dādāe varsādnũ anumān karyũ.}
\end{quote}

\begin{quote}
Looking at the sky, grandfather guessed that rain was coming.
\end{quote}

A neuter noun.

\begin{cousinweb}
\textbf{\gu{અનુમાન}} (\emph{anumān}) is a tatsama, borrowed from Sanskrit \emph{anumāna}, \textquotedblleft{}an inference, a conclusion from signs\textquotedblright{} (Monier-Williams).
\end{cousinweb}

\begin{grammarlens}[title={What you've built}]
One more word for giving reasons and opinions.
\end{grammarlens}

\subsection*{Your turn}
\begin{itemize}
\raggedright
  \item \emph{Say it:} \emph{anumān}
  \item \emph{Say it:} \emph{anumān}, once more
  \item \emph{Say it:} \emph{ākāś joīne dādāe varsādnũ anumān karyũ.} — Looking at the sky, grandfather guessed that rain was coming.
  \item \emph{Recall:} say the Gujarati for a principle, then the Gujarati for everyday practice, then say \emph{anumān} again
\end{itemize}

\begin{tcolorbox}[breakable,colback=teal!4,colframe=teal!35!black,title={Before you move on}]
What does \gu{અનુમાન} mean? (\textquotedblleft{}An inference\textquotedblright{}.) Say it once more.
\end{tcolorbox}
\section[\texorpdfstring{\gu{વાસ્તવિક} (vāstavik)}{વાસ્તવિક (vāstavik)}]{\gu{વાસ્તવિક} (vāstavik) — actual}

\label{lesson:GU-C365-vastavik}



\begin{quote}
Before the new one: say the Gujarati for everyday practice, then the Gujarati for an inference.
\end{quote}

\subsection*{You'll want to know: \gu{વાસ્તવિક}}
\textbf{\gu{વાસ્તવિક}} — \emph{vāstavik} — \textquotedblleft{}actual\textquotedblright{}.

\begin{quote}
\gu{યોજનાની} \gu{વાસ્તવિક} \gu{કિંમત} \gu{ઘણી} \gu{વધારે} \gu{હતી}.
\end{quote}

\begin{quote}
\emph{yojnānī vāstavik kimmat ghaṇī vadhāre hatī.}
\end{quote}

\begin{quote}
The actual cost of the plan was much higher.
\end{quote}

\begin{cousinweb}
\textbf{\gu{વાસ્તવિક}} (\emph{vāstavik}) is a tatsama, borrowed from Sanskrit \emph{vāstavika}, \textquotedblleft{}real, substantial\textquotedblright{} (Monier-Williams).
\end{cousinweb}

\begin{grammarlens}[title={What you've built}]
One more word for giving reasons and opinions.
\end{grammarlens}

\subsection*{Your turn}
\begin{itemize}
\raggedright
  \item \emph{Say it:} \emph{vāstavik}
  \item \emph{Say it:} \emph{vāstavik}, once more
  \item \emph{Say it:} \emph{yojnānī vāstavik kimmat ghaṇī vadhāre hatī.} — The actual cost of the plan was much higher.
  \item \emph{Recall:} say the Gujarati for everyday practice, then the Gujarati for an inference, then say \emph{vāstavik} again
\end{itemize}

\begin{tcolorbox}[breakable,colback=teal!4,colframe=teal!35!black,title={Before you move on}]
What does \gu{વાસ્તવિક} mean? (\textquotedblleft{}Actual\textquotedblright{}.) Say it once more.
\end{tcolorbox}
\section[\texorpdfstring{\gu{સાંકળવું} (sā̃kaḷvũ)}{સાંકળવું (sā̃kaḷvũ)}]{\gu{સાંકળવું} (sā̃kaḷvũ) — to link}

\label{lesson:GU-C365-sakalvu}



\begin{quote}
Before the new one: say the Gujarati for an inference, then the Gujarati for actual.
\end{quote}

\subsection*{You'll want to know: \gu{સાંકળવું}}
\textbf{\gu{સાંકળવું}} — \emph{sā̃kaḷvũ} — \textquotedblleft{}to link\textquotedblright{}.

\begin{quote}
\gu{ગુરુ} \gu{નવો} \gu{પાઠ} \gu{જૂના} \gu{પાઠ} \gu{સાથે} \gu{સાંકળે} \gu{છે}.
\end{quote}

\begin{quote}
\emph{guru navo pāṭh jūnā pāṭh sāthe sā̃kḷe che.}
\end{quote}

\begin{quote}
The teacher links the new lesson to the old one.
\end{quote}

Past, with \gu{તેણે} (he or she …): \gu{તેણે} \gu{બંને} \gu{વિચાર} \gu{સાંકળ્યા}, \emph{teṇe banne vicār sā̃kaḷyā}.

\begin{cousinweb}
\textbf{\gu{સાંકળવું}} (\emph{sā̃kaḷvũ}) is built on \emph{sā̃kaḷ}, \textquotedblleft{}a chain\textquotedblright{}, which CDIAL traces to Sanskrit \emph{śṛṅkhala}, \textquotedblleft{}a chain\textquotedblright{} (CDIAL 12580): to chain two things together.
\end{cousinweb}

\begin{grammarlens}[title={What you've built}]
One more word for giving reasons and opinions.
\end{grammarlens}

\subsection*{Your turn}
\begin{itemize}
\raggedright
  \item \emph{Say it:} \emph{sā̃kaḷvũ}
  \item \emph{Say it:} \emph{sā̃kaḷvũ}, once more
  \item \emph{Say it:} \emph{guru navo pāṭh jūnā pāṭh sāthe sā̃kḷe che.} — The teacher links the new lesson to the old one.
  \item \emph{Recall:} say the Gujarati for an inference, then the Gujarati for actual, then say \emph{sā̃kaḷvũ} again
\end{itemize}

\begin{tcolorbox}[breakable,colback=teal!4,colframe=teal!35!black,title={Before you move on}]
What does \gu{સાંકળવું} mean? (\textquotedblleft{}To link\textquotedblright{}.) Say it once more.
\end{tcolorbox}
